% ==============================================================================
% capitulos/00-referencias.tex - Bibliografía del Proyecto
% ==============================================================================

\begin{thebibliography}{31}

% Paquete opcional para mejor manejo de URLs (ponerlo en el preámbulo)
% \usepackage{xurl}
% \usepackage[hyphens]{url}

\bibitem{1} 
Instituto Politécnico Nacional, \textit{``Ley Orgánica del Instituto Politécnico Nacional''}, Diario Oficial de la Federación, Ciudad de México, México. Disponible en: \url{https://www.ipn.mx/assets/files/defensoria/docs/manuales/ley-org\%C3\%A1nica-ipn.pdf}

\bibitem{2} 
Instituto Politécnico Nacional, \textit{``Reglamento Interno del Instituto Politécnico Nacional''}, Gaceta Politécnica, Ciudad de México, México. Disponible en: \url{https://www.ipn.mx/assets/files/defensoria/docs/manuales/reglamento-interno.pdf}

\bibitem{3} 
Instituto Politécnico Nacional, \textit{``Acuerdo de creación de la Defensoría de los Derechos Politécnicos''}, Gaceta Politécnica, no. 622 (Extraordinaria), pp. 3-7, 31 ene. 2006. Disponible en: \url{https://www.ipn.mx/assets/files/imageninstitucional/docs/gaceta-extraordinaria/2006/01/622-31-enero-gaceta-extra-.pdf}

\bibitem{4} 
Instituto Politécnico Nacional, \textit{``Manual de Organización de la Defensoría de los Derechos Politécnicos''}, Dirección de Planeación y Organización, Ciudad de México, México, oct. 2022. Disponible en: \url{https://www.ipn.mx/assets/files/defensoria/docs/manuales/moddp2022.pdf}

\bibitem{5} 
Instituto Politécnico Nacional, \textit{``Manual de Procedimientos de la Defensoría de los Derechos Politécnicos''}, Dirección de Planeación y Organización, Ciudad de México, México, oct. 2024. Disponible en: \url{https://www.ipn.mx/assets/files/defensoria/docs/manuales/mp-ddp.pdf}

\bibitem{6} 
Secretaría de la Función Pública, \textit{``Sistema Integral de Denuncias Ciudadanas (SIDEC)''}, Gobierno de México. Disponible en: \url{https://sidec.buengobierno.gob.mx/}

\bibitem{7} 
Agencia Mexicana de Cooperación Internacional para el Desarrollo (AMEXCID), \textit{``Catálogo de Capacidades Mexicanas: Sistema Integral de Quejas y Denuncias Ciudadanas (SIDEC)''}, SRE. Disponible en: \url{https://de.sre.gob.mx/capacidades/instituciones-mexicanas/sfp/sistema-integral-de-quejas-y-denuncias-ciudadanas-sidec}

\bibitem{8} 
Secretaría Anticorrupción y Buen Gobierno, \textit{``Implementan SFP, BID y GovLab proyecto de mejora al Sistema Integral de Denuncias Ciudadanas (SIDEC)''}, Gob.mx, 22 jul. 2018. Disponible en: \url{https://www.gob.mx/buengobierno/prensa/implementan-sfp-bid-y-govlab-proyecto-de-mejora-al-sistema-integral-de-denuncias-ciudadanas-sidec}

\bibitem{9} 
Gobierno de la Ciudad de México, \textit{``Denuncia Digital''}, Fiscalía General de Justicia de la CDMX. Disponible en: \url{https://denunciadigital.cdmx.gob.mx/}

\bibitem{10} 
Denuncia.org, \textit{``Denuncia en línea: Consulta si en tu estado la Fiscalía cuenta con la opción de denuncia digital''}, Impunidad Cero, 2026. Disponible en: \url{https://denuncia.org/denuncia-digital/}

\bibitem{11} 
Fiscalía General de Justicia de la Ciudad de México, \textit{``Servicios en Línea: Denuncia Digital, Denuncia Anónima y MP Transparente''}, FGJ-CDMX, 2026. Disponible en: \url{https://www.fgjcdmx.gob.mx/servicios/en-linea}

\bibitem{12} 
Gobierno de la Ciudad de México, \textit{``MP Virtual: Ratificación de Querellas y Actas Especiales''}, FGJ-CDMX. Disponible en: \url{https://mpvirtual.fgjcdmx.gob.mx/CiberDenuncia/TerminosCondiciones.aspx}

\bibitem{13} 
Nación321, \textit{``5 apps que nos ayudan a combatir la corrupción''}, Grupo Multimedia Lauman, SAPI de CV, 16 feb. 2024. Disponible en: \url{https://www.nacion321.com/ciudadanos/5-apps-que-nos-ayudan-a-combatir-la-corrupcion}

\bibitem{14} 
Secretaría de la Contraloría del Gobierno del Estado de México, \textit{``Sistema de Atención Mexiquense (SAM)''}, GEM, 2026. Disponible en: \url{https://www.secogem.gob.mx/sam/tramite.asp}

\bibitem{15} 
Whistlelink, \textit{``The All-in-One Whistleblowing Solution''}, Whistlelink.com, 2026. Disponible en: \url{https://www.whistlelink.com/}

\bibitem{16} 
EQS Group, \textit{``EQS Integrity Line: The leading whistleblowing system in Europe''}, EQS.com, 2026. Disponible en: \url{https://www.eqs.com/es/soluciones-gobierno-corporativo/integrity-line/}

\bibitem{17} 
NAVEX Global, \textit{``WhistleB: Whistleblowing Solutions for a Modern Workplace''}, Navex.com, 2026. Disponible en: \url{https://www.navex.com/en-us/products/whistleb/}

\bibitem{18} 
Whistleblower Software by Whistleblowing Solutions ApS, \textit{``The simple, secure, and compliant whistleblowing system''}, Whistleblowersoftware.com, 2026. Disponible en: \url{https://whistleblowersoftware.com/es}

\bibitem{19} 
Lefebvre, \textit{``Centinela: Canal de denuncias para el cumplimiento normativo''}, Lefebvre.es, 2026. Disponible en: \url{https://lefebvre.es/soluciones/centinela/canal-denuncias/}

\bibitem{20} 
Coloriuris, \textit{``Canal de Denuncias: Cumplimiento de la Ley 2/2023 y Directiva (UE) 2019/1937''}, Coloriuris.net, 2026. Disponible en: \url{https://www.coloriuris.net/canal-de-denuncias/}

\bibitem{21} 
Bully Button, \textit{``The Bully Button: Reporting and Prevention System''}, BullyButton.com, 2026. Disponible en: \url{http://bullybutton.com/}

\bibitem{22} 
Secretaría de Educación Pública, \textit{``Ventanilla Escolar: Orientación y Quejas''}, Gobierno de México. Disponible en: \url{https://www.gob.mx/sep/acciones-y-programas/ventanilla-escolar}

\bibitem{23} 
Instituto Politécnico Nacional, \textit{``Protocolo Alerta IPN: Seguridad Institucional''}, Secretaría General del IPN. Disponible en: \url{https://www.ipn.mx/seguridad/}

\bibitem{24} 
Universidad Nacional Autónoma de México, \textit{``Sistema de Denuncia en Línea de la Defensoría de los Derechos Universitarios''}, UNAM. Disponible en: \url{https://www.defensoria.unam.mx/denuncia/}

\bibitem{25} 
Procuraduría Federal de la Defensa del Trabajo, \textit{``PROFEDET Digital: Servicios de orientación y seguimiento''}, Gobierno de México. Disponible en: \url{https://www.gob.mx/profedet}

\bibitem{26} 
Instituto Politécnico Nacional, \textit{``Protocolo para la Prevención, Detección, Atención y Sanción de la Violencia de Género en el Instituto Politécnico Nacional''}, Gaceta Politécnica, Número 1726, Ciudad de México, México, mayo 2023.

\bibitem{27} 
BBVA, \textit{``Arquitectura de software o sistemas: ¿qué es y ejemplos?''}, bbva.com, 2021. [En línea]. Disponible: \url{https://www.bbva.com/es/innovacion/arquitectura-de-software-o-sistemas-que-es-y-ejemplos/}. [Accedido: 10-abr-2026].

\bibitem{28} 
UTEL, \textit{``Atributos de calidad del software''}, SCALA. [En línea]. Disponible: \url{https://gc.scalahed.com/recursos/files/r161r/w26023w/L1IDS118_l2_s3.pdf}. [Accedido: 10-abr-2026].

\bibitem{29} 
Microsoft, \textit{``¿Qué es la escalabilidad?''}, Microsoft Learn, 2023. [En línea]. Disponible: \url{https://learn.microsoft.com/es-es/biztalk/core/what-is-scalability}. [Accedido: 10-abr-2026].

\bibitem{30} 
IBM, \textit{``¿Qué es una arquitectura monolítica?''}, ibm.com. [En línea]. Disponible: \url{https://www.ibm.com/think/topics/monolithic-architecture}. [Accedido: 10-abr-2026].

\bibitem{31} 
Amazon Web Services (AWS), \textit{``Construcción de una arquitectura distribuida y escalable''}, aws.amazon.com. [En línea]. Disponible: \url{https://aws.amazon.com/startups/learn/building-scalable-distributed-architecture?lang=es}. [Accedido: 10-abr-2026].

\bibitem{32} 
Amazon Web Services (AWS), "¿Qué son los microservicios?," aws.amazon.com. [En línea]. Disponible: https://aws.amazon.com/es/microservices/. [Accedido: 10-abr-2026].

\bibitem{33} 
IBM, \textit{``¿Qué es la arquitectura de tres capas?''}, ibm.com. [En línea]. Disponible: \url{https://www.ibm.com/mx-es/think/topics/three-tier-architecture}. [Accedido: 10-abr-2026].

\bibitem{34} 
NGINX, \textit{``nginx''}, nginx.org. [En línea]. Disponible: \url{https://nginx.org/en/}. [Accedido: 10-abr-2026].

\bibitem{35} 
B. Khanal, \textit{``Nginx: The Lightweight Powerhouse for Modern Web Serving''}, Medium, 2023. [En línea]. Disponible: \url{https://medium.com/@bishal.khanal.2057/nginx-the-lightweight-powerhouse-for-modern-web-serving-717f7cc628da}. [Accedido: 10-abr-2026].

\bibitem{36}
Oracle, \textit{``¿Qué es la tecnología Java y para qué la necesito?''}, java.com. [En línea]. Disponible: \url{https://www.java.com/es/download/help/whatis_java.html}. [Accedido: 10-abr-2026].

\bibitem{37}
IBM, \textit{``¿Qué es Java?''}, ibm.com. [En línea]. Disponible: \url{https://www.ibm.com/mx-es/think/topics/java}. [Accedido: 10-abr-2026].

\bibitem{38}
IBM, \textit{``¿Qué es Spring Boot?''}, ibm.com. [En línea]. Disponible: \url{https://www.ibm.com/mx-es/think/topics/java-spring-boot}. [Accedido: 10-abr-2026].

\bibitem{39}
Amazon Web Services (AWS), \textit{``¿Qué es Python?''}, aws.amazon.com. [En línea]. Disponible: \url{https://aws.amazon.com/es/what-is/python/}. [Accedido: 10-abr-2026].

\bibitem{40}
Angular, \textit{``What is Angular?''}, angular.dev. [En línea]. Disponible: \url{https://angular.dev/overview}. [Accedido: 10-abr-2026].

\bibitem{41}
Imagina Formación, \textit{``Conoce las ventajas de utilizar Angular''}, imaginaformacion.com. [En línea]. Disponible: \url{https://imaginaformacion.com/tutoriales/conoce-las-ventajas-de-utilizar-angular}. [Accedido: 10-abr-2026].

\bibitem{42}
IBM, \textit{``¿Qué es PostgreSQL?''}, ibm.com. [En línea]. Disponible: \url{https://www.ibm.com/mx-es/think/topics/postgresql}. [Accedido: 10-abr-2026].

\bibitem{43}
Red Hat, \textit{``¿Qué es Docker?''}, redhat.com. [En línea]. Disponible: \url{https://www.redhat.com/es/topics/containers/what-is-docker}. [Accedido: 10-abr-2026].

\bibitem{44}
GitHub, \textit{``Acerca de GitHub y Git''}, docs.github.com. [En línea]. Disponible: \url{https://docs.github.com/es/get-started/start-your-journey/about-github-and-git}. [Accedido: 10-abr-2026].

\bibitem{45}
E. Linis, \textit{``Conceptos Básicos de Git''}, Gist GitHub. [En línea]. Disponible: \url{https://gist.github.com/erlinis/57a55dfb0337f5cd15cd}. [Accedido: 10-abr-2026].

\bibitem{46}
Villamizar, M., et al. \textit{Evaluating the Monolithic and the Microservice Architecture Pattern to Deploy Web Applications in the Cloud}. Universidad de los Andes, Colombia.
\bibitem{47}
Tanenbaum, A. S., et al. \textit{Arquitectura de Sistemas Distribuidos}.

\bibitem{48} 
Microsoft Ibérica. \textit{Guía de Arquitectura N-Capas orientada al Dominio con .NET 4.0}.

\bibitem{49}
Poza Luján, J. L. \textit{Propuesta de arquitectura distribuida de control inteligente basada en políticas de calidad de servicio}. Universitat Politècnica de València.

\bibitem{50}
Coti Colop, B. M. \textit{Reglas del Negocio en Arquitectura Tres Capas}. Universidad de San Carlos de Guatemala.

\bibitem{51}
Álvarez Caules, C. \textit{Spring Boot Actuator}. Arquitectura Java. Recuperado de: https://www.arquitecturajava.com/spring-boot-actuator/

\bibitem{52}
Spring Security. \textit{OAuth 2.0 Resource Server JWT}. VMware, Inc. Recuperado de: https://docs.spring.io/spring-security/reference/servlet/oauth2/resource-server/jwt.html

\bibitem{53}
Spring Cloud. \textit{Spring Cloud Config \& Netflix Eureka}. VMware, Inc. Recuperado de: https://spring.io/projects/spring-cloud

\bibitem{54}
Fielding, R. T. \textit{Architectural Styles and the Design of Network-based Software Architectures}. University of California, Irvine.

\bibitem{55}
Kleppmann, M. \textit{Designing Data-Intensive Applications: The Big Ideas Behind Reliable, Scalable, and Maintainable Systems}. O'Reilly Media.

\bibitem{56}
Rescorla, E. \textit{The Transport Layer Security (TLS) Protocol Version 1.3}. RFC 8446, Internet Engineering Task Force (IETF).

\bibitem{57}
Stopford, B. \textit{Designing Event-Driven Systems: Concepts and Patterns for Streaming Services with Apache Kafka}. O'Reilly Media.

\bibitem{58}
Nygard, M. T. \textit{Release It!: Design and Deploy Production-Ready Software}. Pragmatic Bookshelf.

\bibitem{59}
Burns, B., Beda, J., y Hightower, K. \textit{Kubernetes: Up and Running: Dive into the Future of Infrastructure}. O'Reilly Media.

\bibitem{60}
Schönig, H. J. \textit{Mastering PostgreSQL 15: Advanced techniques to build and manage scalable, reliable, and fault-tolerant database applications}. Packt Publishing.

\bibitem{cortes1995}
C. Cortes y V. Vapnik, \textit{``Support-Vector Networks''}, \textit{Machine Learning}, vol. 20, no. 3, pp. 273--297, sep. 1995. DOI: \url{https://doi.org/10.1007/BF00994018}
 
\bibitem{breiman2001}
L. Breiman, \textit{``Random Forests''}, \textit{Machine Learning}, vol. 45, no. 1, pp. 5--32, oct. 2001. DOI: \url{https://doi.org/10.1023/A:1010933404324}
 
\bibitem{vaswani2017}
A. Vaswani, N. Shazeer, N. Parmar, J. Uszkoreit, L. Jones, A. N. Gomez, Ł. Kaiser e I. Polosukhin, \textit{``Attention is All You Need''}, en \textit{Advances in Neural Information Processing Systems 30 (NIPS 2017)}, pp. 5998--6008, 2017. Disponible en: \url{https://arxiv.org/abs/1706.03762}
 
\bibitem{reimers2019}
N. Reimers e I. Gurevych, \textit{``Sentence-BERT: Sentence Embeddings using Siamese BERT-Networks''}, en \textit{Proceedings of the 2019 Conference on Empirical Methods in Natural Language Processing (EMNLP-IJCNLP)}, pp. 3982--3992, Hong Kong, China, nov. 2019. Disponible en: \url{https://arxiv.org/abs/1908.10084}
 
\bibitem{mikolov2013}
T. Mikolov, I. Sutskever, K. Chen, G. S. Corrado y J. Dean, \textit{``Distributed Representations of Words and Phrases and their Compositionality''}, en \textit{Advances in Neural Information Processing Systems 26 (NIPS 2013)}, pp. 3111--3119, 2013. Disponible en: \url{https://arxiv.org/abs/1310.4546}
 
\bibitem{honnibal2020}
M. Honnibal, I. Montani, S. Van Landeghem y A. Boyd, \textit{``spaCy: Industrial-strength Natural Language Processing in Python''}, Zenodo, 2020. DOI: \url{https://doi.org/10.5281/zenodo.1212303}
 
\bibitem{pedregosa2011}
F. Pedregosa \textit{et al.}, \textit{``Scikit-learn: Machine Learning in Python''}, \textit{Journal of Machine Learning Research}, vol. 12, pp. 2825--2830, 2011. Disponible en: \url{http://www.jmlr.org/papers/volume12/pedregosa11a/pedregosa11a.pdf}


\end{thebibliography}